\section*{September 8, 2026: Initial verdicts on remaining land and locations}
\textit{Agent-drafted review. Jacob requested an initial verdict for every remaining property. These recommendations have not been applied as new production decisions.}

The refreshed list covers 140 source project records and 146 review questions,
including all 80 currently unresolved construction-year geometries. These are
not 140 distinct physical buildings or newly discovered errors. The original
30 parcel requests remain represented: ten now resolve geometrically and twenty
still require attention. Closing a map question does not independently establish
that a building's units and land refer to the same object.

The most promising correction is to recognize obsolete parcel numbers before
trying to geocode them as additional construction. Twenty-one homes at
3613/3615/3625 South Morgan, four at 3763 South Morgan, and ten at 4647/4649
South Lake Park have plausible individually retained successors. All 35 old
recorded points have exactly one current eligible one-home candidate within
one foot; the maximum paired difference is 0.986 feet. The old assessments end
immediately before the successor assessments begin, and land areas agree
exactly or nearly exactly. Construction-year and floor-area revisions remain
explicit evidence to reconcile, not additional homes. The initial recommendation
is to retain the successor homes once after verifying the general parcel crosswalk.
At 3535/3541 South Maplewood, ten individually retained homes each report
2,060 building square feet, exactly matching the old ten-card aggregate's
20,600 square feet. The common-land difference still needs verification.

At 6026 South Green, the original 2014 assessment supplies unusually strong
evidence of a mistyped parcel reference. PIN 20174140310000 reports six dwellings,
9,036 building square feet, 3,100 land square feet and a 60 percent share.
PIN 20174140320000 reports the same building fields, 3,718 land square feet
and a 40 percent share, but points to 29174140310000. That incorrect reference
joins land roughly 12.6 miles away. The initial recommendation is to correct
the link, verify the actual two construction-year polygons and retain one
six-unit building. No link or density value was changed in this review.

Three individually recorded homes at 2504 South King, 2502 South King and
2501 South Calumet receive a shared development-center distance of 463 feet.
Their exact current parcel points are approximately 597, 616 and 615 feet from
the construction-year ward boundary. These are comparison locations, not
newly adopted historical coordinates. They demonstrate why a large former
development center cannot automatically locate each individual home.
The recommendation is to preserve individual measurements and establish
individual historical locations before assigning sample inclusion.

Other specific concerns are recorded property by property in the audit memo.
They include a likely parcel-number replacement at 1373 East 64th; a Morgan
group that combines an unrelated nineteenth-century building with duplicated
townhouse cards; individual Deming homes with repeated whole-development land;
and commercial-source holdings that combine several buildings or phases.
Five late-appearing Vernon parcel numbers warrant a successor-identity check
against the approved five-home site. That potential overlap does not overturn
the settled 2020 site measurement decision.

The memo distinguishes direct source checks from earlier review notes.
The older notes are useful evidence leads, but some propose grouping townhomes
contrary to Jacob's newer individual-building preference. They were not copied
back into production. Seven records have sufficient parcel-area and location
agreement to recommend closing their specific land/location flags; this is
not independent certification of their construction dates or a release approval.

\textit{Validation:} Both new audit outputs build through Make, have standard
data reports, and are up to date on an unchanged build. The location audit asserts
coverage of every unresolved geography record and unique project keys. The
successor audit reports ambiguity rather than treating the nearest point as
verified identity. All 35 proposed paired points match the independently
examined successor records uniquely within one foot. The existing exhaustive
boundary-distance check was rebuilt and passed. The full memo contains one
initial verdict for each of the 140 source project records.

No production cleaning rules, exclusions, construction years or coordinates were
changed. The paper's frozen input remains unchanged, and density regressions
were not rerun. Neither set of 35 Morgan/Lake Park source or successor identifiers
appears under those identifiers in the frozen paper input. The Maplewood aggregate
and its ten individual identifiers do appear, all outside 500 feet, as do both
64th Street identifiers. These are identifier checks, not a complete estimate
of the sample or regression impact. Earlier unrelated logbook outputs were held at their existing
versions solely to compile and inspect this entry.
